\section{Introduction}\label{sec:intro}

Haiden, Katzarkov, Kontsevich and Pandit \cite{hkkp26} propose a notion of
K\"ahler geometry on a triangulated category with a stability condition in
the sense of Bridgeland \cite{bri07}. In their framework an object carries a
space of metrics $h$ and a complex-valued function $\cS_{\CC}$ on it with
differential $\delta h\mapsto\tr(\Omega\,h^{-1}\delta h)$, for an endomorphism
$\Omega$ whose trace is the central charge of the object. A \emph{phase}
$\arg\Omega$ is defined by a positivity condition, the \emph{mass} is
$m=\tr|\Omega|$, and the metrics move by the flow
$h^{-1}\dot h=-\arg\Omega$, called the minimizing flow in
\cite[\S4.1]{hkkp26}. In \cite[\S\S1,~4.1]{hkkp26} the mass is expected to
decrease along the flow, the flow to converge to a \emph{harmonic} metric,
one of constant phase, when the object admits one, and harmonic metrics to
exist on polystable objects, a categorical form of the theorem of Donaldson,
Uhlenbeck and Yau \cite{don85,uy86}. For a representation of a quiver,
$\Omega$ is the central charge of each vertex shifted by the moment map, and
the harmonic metrics are the zeros of the shifted moment map. For complexes
of quiver representations $\Omega$ is defined through a superconnection in a
graded algebra \cite[\S5.4]{hkkp26}, and it is asked there whether the mass
decreases \cite[Question~5.18]{hkkp26}.

For a complex the endomorphism $\Omega$ is in general not normal
\cite[Example~2.3]{wue26}, whereas in \cite[\S5.4]{hkkp26} its phase is
written with the functional calculus. Throughout this paper $\arg\Omega$ is
taken in the sense of the positivity condition of \cite[\S4.2]{hkkp26}, that
is, as the argument of the unitary factor in the polar decomposition
(Section~\ref{sec:setting}). Every statement about the flow below, including
the comparison with \cite[Question~5.18]{hkkp26}, refers to this reading. We
call the resulting flow the \emph{phase flow}: it is the minimizing flow of
\cite[\S4.1]{hkkp26} with this reading of the phase.

The simplest complexes on which both a differential and the arrows of a
quiver act are
\begin{equation}\label{eq:introcomplex}
 (V\xrightarrow{\,(A_j)\,}W)\ \xrightarrow{\,(D,0)\,}\ (U\to0),
\end{equation}
two-term complexes of representations of the generalized Kronecker quiver
with $k$ arrows, with $D:V\to U$ surjective and $A_1,\dots,A_k:V\to W$. The
cohomology is the representation $\ker D\to W$ in degree zero; it is a
subcomplex, and the quotient $(V/\ker D\to0)\to(U\to0)$ is an acyclic complex
of rank $r=\dim U$, the \emph{acyclic part}. So \eqref{eq:introcomplex} is
quasi-isomorphic to a representation. For these complexes $\Omega$ was
computed in \cite{wue26}, where it was shown that the mass can increase along
the phase flow \cite[Theorem~1.1]{wue26}; that the contribution of the
acyclic part to $\Omega$ nevertheless decays if the least singular
value of $D$ in the initial metrics is large \cite[Theorem~1.3]{wue26}; and
that for one pair of charges and $\dim\ker D=\dim W$ the moment map of the
cohomology representation then tends to zero if its Kempf--Ness function is
bounded below \cite[Theorem~1.4]{wue26}.

We prove that, for every initial metric, the contribution of the acyclic
part decays, for all charges, and that, when the charge at the source vertex
has the larger argument, the metrics induced on the cohomology converge,
after a scalar normalization, to a harmonic metric of the cohomology
representation whenever that representation is polystable. This
removes the assumption on the initial metrics of
\cite[Theorems~1.3,~1.4]{wue26}. The main results are
Theorems~\ref{thm:introA} and~\ref{thm:introB}.

\subsection{Results}
Let $h=(h_V,h_U,h_W)$ be Hermitian metrics on the three spaces of
\eqref{eq:introcomplex}; adjoints, written ${}^\dagger$, are taken for the
current metrics. Let $z_0,z_1$ be central charges in the upper half-plane,
attached to the source vertex $0$, which carries $V$ and $U$, and to the
target vertex $1$, which carries $W$; let $n_0=\dim\ker D$, $n_1=\dim W$,
$Z=n_0z_0+n_1z_1$ and $\varphi=\arg Z\in(0,\pi)$. The pair $(z_0,z_1)$ is the
central charge of a stability condition on the bounded derived category of
representations whose heart is the category of representations; the phase of
a representation of dimension vector $(a,b)$ is $\arg(az_0+bz_1)$, which is
$\pi$ times the usual normalized phase. The endomorphism of
\cite[\S5.4]{hkkp26} is $\Omega=\Omega_V\oplus\Omega_U\oplus\Omega_W$; it is
computed in \eqref{eq:coeff} below. Put
\[
 L=DD^\dagger,\qquad\lambda=\lambda_{\min}(L),\qquad
 D^+=D^\dagger(DD^\dagger)^{-1},\qquad s=\sum_j\norm{A_jD^+}_{\HS}^2 .
\]
For the metrics induced on $\ker D$ and $W$ let $H_j$ be the matrix of
$A_j|_{\ker D}$ in orthonormal frames, $C_0=\sum_jH_j^\dagger H_j$ and
$C_1=\sum_jH_jH_j^\dagger$. For the cohomology representation with these
metrics, $\Omega$ is $(z_0+C_0)\oplus(z_1-C_1)$, a scalar standing for the
corresponding multiple of the identity. We write $\norm\cdot_1$ for
the trace norm, and $(z_0+C_0)\oplus0$ for the endomorphism of
$V=\ker D\oplus(\ker D)^\perp$ that is $z_0+C_0$ on $\ker D$ and zero on the
orthogonal complement.

\begin{theorem}[Cancellation of the acyclic part]\label{thm:introA}
For all charges $z_0,z_1$ in the upper half-plane and every complex
\eqref{eq:introcomplex}, along the phase flow $\lambda(t)\to\infty$,
$s(t)\to0$, and for $t\geq t_1$
\[
 \norm{\Omega_V-(z_0+C_0)\oplus0}_1+\norm{\Omega_U}_1+\norm{\Omega_W-(z_1-C_1)}_1
 \leq C_{\mathrm{can}}\,e^{-\frac14(\arg z_1)(t-t_1)},
\]
where $C_{\mathrm{can}}$ is an explicit function of the dimensions and the
charges, and $t_1$ of these and of three numbers computed from the initial
metrics. No hypothesis is made on the cohomology representation.
\end{theorem}

The charges are \emph{ordered} if $\arg z_0>\arg z_1$, that is, if the simple
representation at the source vertex has the larger phase, and
\emph{reversed} if $\arg z_0<\arg z_1$. For reversed charges no
representation with $n_0,n_1\geq1$ is semistable. For ordered charges there
are unique $\theta_0,\theta_1>0$ with
$\arg(z_0+\theta_0)=\arg(z_1-\theta_1)=\varphi$; then
$n_0\theta_0=n_1\theta_1$, and a representation of dimension vector
$(n_0,n_1)$ is semistable, in the sense that no subrepresentation has phase
larger than $\varphi$, exactly when it is semistable in the sense of King
\cite{kin94} for the character $(-\theta_0,\theta_1)$. We call
$(\theta_0,\theta_1)$ the \emph{stability parameter}. The \emph{moment map} of
the cohomology representation, shifted by the stability parameter, is
$\mu=(C_0-\theta_0,\,C_1-\theta_1)$. Its zeros are the pairs of metrics for
which the phases $\arg(z_0+C_0)$ and $\arg(z_1-C_1)$ at the two vertices both
equal $\varphi$, that is, the harmonic metrics; they exist exactly for polystable
representations \cite{kn79,kin94}.

\begin{theorem}[Convergence]\label{thm:introB}
Along the phase flow, for every initial metric:
\begin{enumerate}
\item[(i)] If the charges are ordered and the cohomology representation is
 semistable, then $\mu(t)\to0$ and $m(h(t))\to|Z|$.
\item[(ii)] If the charges are ordered and the cohomology representation is
 polystable, the induced metrics
 $(e^{\varphi t}h_V|_{\ker D},\,e^{\varphi t}h_W)$ converge exponentially, as
 Hermitian forms on the fixed vector spaces $\ker D$ and $W$, to a zero of
 the moment map, that is, to a harmonic metric of the cohomology
 representation; the moment map and $m-|Z|$ tend to zero exponentially.
\item[(iii)] If the charges are reversed, the arrows $H_j$ tend to zero
 exponentially and $m(h(t))\to n_0|z_0|+n_1|z_1|$.
\end{enumerate}
\end{theorem}

A harmonic metric of a representation is not stationary under the flow but
moves by the scalar $e^{-\varphi t}$, which is why the limit is stated for
the normalized metrics. The complex itself has no harmonic metric, since the
phase on $U$ is negative and the phases on $V,W$ are positive, and its
metrics do not converge even after normalization, because $\lambda\to\infty$;
the limit exists for the metrics induced on the cohomology. In
\cite[Remark~5.15]{hkkp26} the metrics of an object are the metrized
complexes isomorphic to it in the derived category. The flow preserves the
underlying complex, and Theorem~\ref{thm:introB} identifies its limit with a
harmonic metric on another representative of the same object, the
cohomology representation. The rate in (ii)
is governed by the least eigenvalue of the Hessian of the Kempf--Ness
function transversally to the zeros of the moment map and by the rate of
Theorem~\ref{thm:introA}. The precise statements are
Theorems~\ref{thm:cancellation}, \ref{thm:relaxation}, \ref{thm:convergence}
and~\ref{thm:reversed}.

The proof of Theorem~\ref{thm:introA} does not use a Lyapunov function: the
mass is not monotone, and the function $\cK$ introduced before
Theorem~\ref{thm:introE} is monotone for the charges of that theorem but not
for all charges. It rests on an a priori bound for the
following quantity, which is not monotone either, and on the growth of
$\lambda$. Let $R(b)=(1-e^{-b})/b$ and
\[
 N=\sum_j\tr\bigl(A_jR(D^\dagger D)A_j^\dagger\bigr),
\]
the \emph{arrow energy}; it is the term of $-\cS_{\CC}$, see
\eqref{eq:primitive}, that does not contain the charges, and it bounds $\tr C_0$ and $(1-e^{-\lambda})s$.

\begin{theorem}[A priori bound]\label{thm:introC}
For all charges in the upper half-plane, along the phase flow
\[
 \frac{dN}{dt}\leq-\rho N+n_1M_1,\qquad
 \rho=\pi-\max\bigl\{\arg z_0,\tfrac{2\pi}5\bigr\}>0,
\]
with $M_1=\max\{\pi\max(\Rea z_1,0),\,2\Ima z_1\}$. Hence
$N(t)\leq\max\{N(0),n_1M_1/\rho\}$ for all $t\geq0$, and
$N(t)\leq2n_1M_1/\rho$ after an explicit time.
\end{theorem}

The three terms of $dN/dt$ come from the three blocks of $\Omega$
(Lemma~\ref{lem:energy}). The block $\Omega_W$ is a normal matrix and its
term is elementary. The blocks $\Omega_V$ and $\Omega_U$ are not normal, and
their terms are controlled by two matrix inequalities, stated and proved
independently of the flow in Sections~\ref{sec:kernel} and~\ref{sec:phase}.

\begin{theorem}[Two matrix inequalities]\label{thm:introD}
The following hold.
\begin{enumerate}
\item[(i)] Let $L>0$ and $E\geq0$ be Hermitian matrices, and let $J_L(E)$ be
 the integral of $e^{-xL}Ee^{-yL}$ over the simplex $x,y\geq0$, $x+y\leq1$.
 Then the negative part of the Hermitian matrix $J_L(E)$ satisfies
 $\tr J_L(E)_-\leq\frac25\tr J_L(E)$.
\item[(ii)] Let $\Omega',\Omega''$ be matrices with the same positive
 definite imaginary part and $\Phi',\Phi''$ the arguments of the unitary
 factors in their polar decompositions, with spectrum in $(0,\pi)$. Then
 $\tr\bigl((\Omega''-\Omega')(\Phi''-\Phi')\bigr)\leq0$, with equality only
 if $\Omega'=\Omega''$.
\end{enumerate}
\end{theorem}

By the Hermite--Genocchi formula for divided differences, in an eigenbasis
of $L$ the matrix
$J_L(E)$ is the entrywise product of $E$ with the matrix of second divided
differences of $e^{-x}$ at the points $\ell_i,0,\ell_k$; this matrix is not
positive semidefinite, and (i) says that the negative part of the product
carries at most two fifths of the trace. The best constant is at least
$(\sqrt2-1)/2$. The proof reduces (i) to an inequality in one real variable,
which is proved by estimates near $0$ and $\infty$ and by interval
arithmetic on a compact interval (Section~\ref{sec:kernel}). Statement (ii)
says that minus the argument of the polar factor is a strictly monotone
operator, in the sense of Minty \cite{min62}, on the Hermitian matrices with
the trace inner product; it is not monotone for the L\"owner order, because
the
arctangent is not operator monotone on the real line
(Remark~\ref{rem:phasefails}). For a scalar imaginary part it is the fact
that the argument of a complex number decreases in its real part. The proof
uses an integral representation of anti-L\"owner type \cite{aud11}
(Section~\ref{sec:phase}).

The mass is not monotone along the phase flow, also for the charges and
complexes of Theorem~\ref{thm:introB}: we give a complex with $512$ arrows,
stable cohomology and $r=1$, with metrics at which the mass increases, for
$z_0=-1+\ii$, $z_1=1+\ii$ (Proposition~\ref{prop:mass}). The largest phase
can increase and the smallest phase can decrease
(Proposition~\ref{prop:extreme}), whereas for representations the mass and
the extreme phases are monotone, as expected in \cite[\S4.1]{hkkp26}
(Proposition~\ref{prop:representations}). What converges for the complexes
\eqref{eq:introcomplex}, for ordered charges and semistable cohomology, is
the mass measure of \cite[\S4.2]{hkkp26}, recalled in
Section~\ref{sec:expectations} (Proposition~\ref{prop:measure}).

For a range of charges there is a Lyapunov function. The \emph{phase
potential} is
\[
 \cK=\frac{\Ima(e^{-\ii\varphi}\cS_{\CC})}{\sin\varphi},
\]
with $\cS_{\CC}$ the function \eqref{eq:primitive}. We call
$z_0=-1+\ii$, $z_1=1+\ii$ the \emph{symmetric charges}; for them and
$n_0=n_1$,
\[
 \cK(h)=\log\det h_V-\log\det h_U-\log\det h_W+\tr\Ein(D^\dagger D)+N,
 \qquad\Ein(b)=\int_0^1(1-e^{-xb})\frac{dx}x ,
\]
with the determinants of the Gram matrices in fixed reference bases.

\begin{theorem}[Lyapunov function]\label{thm:introE}
Let the charges be ordered, and suppose that $\varphi\geq\frac{2\pi}5$ or
$r=1$. Then the phase potential is strictly decreasing along the phase flow.
For the symmetric charges and $n_0=n_1$,
\[
 -\frac{d\cK}{dt}\ \geq\ \frac\pi{10}\tr J_L(DCD^\dagger)+\frac\pi2\tr e^{-L}>0,
 \qquad C=\sum_jA_j^\dagger A_j .
\]
There are ordered charges and complexes with $r=2$, with metrics at which
the phase potential increases along the flow.
\end{theorem}

The number $\frac{2\pi}5$ is $\pi$ times the constant of
Theorem~\ref{thm:introD}(i). The rate of the potential is a sum of three
terms, one for each of $V$, $W$ (degree zero) and $U$ (degree one); the first
two are nonnegative for elementary reasons, and the third is controlled by
Theorem~\ref{thm:introD}(i). See Theorem~\ref{thm:monotone} and
Propositions~\ref{prop:reversed}, \ref{prop:orderedfailure}.

\begin{theorem}[Infimum and moduli]\label{thm:introF}
Let the charges be ordered, and let $\gamma$ be Euler's constant.
\begin{enumerate}
\item[(i)] The potential depends on reference bases of $V,U,W$ only through
 their volume elements. With volume elements on $\ker D$ and $W$ induced
 from those of $V,U,W$ through $\det V\cong\det\ker D\otimes\det U$, the
 infimum of the phase potential over all metrics of the complex is the
 infimum of the Kempf--Ness function \eqref{eq:KN} of its cohomology
 representation for the stability parameter $(\theta_0,\theta_1)$, plus
 $\gamma\,\theta_0r$; it is not attained. Consequently
 the infima for two quasi-isomorphic complexes of the class
 \eqref{eq:introcomplex} with semistable cohomology and corresponding volume
 elements differ by $\gamma\,\theta_0(r'-r)$.
\item[(ii)] On the stable locus of a holomorphic family of complexes,
 $\ii\partial\bar\partial$ of the infimum is the pullback of the K\"ahler form
 of the moduli space of stable representations.
\item[(iii)] For semistable cohomology the potential along the flow tends to
 its infimum, exponentially and locally uniformly on the stable locus.
\end{enumerate}
\end{theorem}

Part (ii) follows from (i) and the classical description of the K\"ahler
potential of a K\"ahler quotient (Lemma~\ref{lem:kappa}). See
Theorems~\ref{thm:infimum}, \ref{thm:moduli} and~\ref{thm:families}.

\subsection{Remarks on the setting and the method}
The phase flow is not the gradient flow of the phase potential
(Section~\ref{sec:potential}). For a representation
regarded as a complex in degree zero $\Omega$ is normal, the phase
potential is the Kempf--Ness function, and its monotonicity along the flow
follows from formulas displayed in \cite[\S\S4.2, 5.1]{hkkp26}. On a
representation the flow has the same stationary points modulo scaling as the
downward gradient flow of the Kempf--Ness function, and its velocity pairs
negatively with the gradient, but it is not that gradient flow
(Section~\ref{sec:convergence}).

The estimates for the polar decomposition that isolate one block of a
matrix, and the differential inequality for $s$, are taken from \cite{wue26}
and recalled in Proposition~\ref{prop:imported}. What is new in the proof of
Theorem~\ref{thm:introA} is the control of the arrows and of $\lambda$ before
$\lambda$ is large: Theorem~\ref{thm:introC} bounds the arrows, and a second
differential inequality, for the trace of the exponential integral of $L$,
forces $\lambda$ to grow (Theorem~\ref{thm:growth}). For the charges of
Theorem~\ref{thm:introE} and semistable cohomology the Lyapunov function
gives the same control with different constants
(Proposition~\ref{prop:apriori}). Theorem~\ref{thm:introB} then follows from
the convexity of the Kempf--Ness function; the relevant statements are
classical \cite{kn79,grs21} and are proved for the present linear action in
Section~\ref{sec:convergence}.

The following is not proved here. For semistable cohomology that is not
polystable the induced metrics diverge
(Proposition~\ref{prop:divergence}); the asymptotics are determined here
only for the smallest such representation, regarded as a complex in degree
zero, where the metrics have the form $t^{\mp1/2}$ on a subrepresentation
and the quotient (Proposition~\ref{prop:semistable}). In general the
normalized induced metrics stay at bounded distance from a trajectory of the
gradient flow of the Kempf--Ness function for the invariant Riemannian
metric weighted by the imaginary parts of the charges
\cite[Theorem~6.3 and Corollary~6.5]{wue26b}, whose asymptotics are
described in \cite{hkkp23}; whether the bounded factor converges is open
beyond this example. The case
$\arg z_0=\arg z_1$ is covered by Theorems~\ref{thm:introA}
and~\ref{thm:introC} but not by Theorem~\ref{thm:introB}. Complexes whose
term of degree one sits at the target vertex, with acyclic parts at both
vertices, or with terms in more than two degrees, are not treated. For
cohomology that is not semistable, Theorems~\ref{thm:introA}
and~\ref{thm:introC} apply, and the induced metrics follow the
Harder--Narasimhan filtration of the cohomology representation
\cite[Theorem~6.11]{wue26b} (Remark~\ref{rem:unstable}).

\subsection{Organization}
Section~\ref{sec:setting} recalls the endomorphism $\Omega$, the flow and the
estimates
used from \cite{wue26}. Sections~\ref{sec:kernel} and~\ref{sec:phase} prove
Theorem~\ref{thm:introD}. Section~\ref{sec:energy} proves
Theorems~\ref{thm:introC} and~\ref{thm:introA}, and
Section~\ref{sec:convergence} Theorem~\ref{thm:introB}.
Section~\ref{sec:potential} treats the phase potential and proves
Theorem~\ref{thm:introE}, and Section~\ref{sec:moduli} proves
Theorem~\ref{thm:introF}. Section~\ref{sec:expectations} treats the mass, the
extreme phases and the fixed points, and Section~\ref{sec:remarks} the charges
for which the phase potential is not monotone, and open questions.
